% Reporte de la Actividad 9
% Semaforo binario y contador
\documentclass[12pt,oneside]{template/liverpoolthesis}
% !TEX root = ../Practica1_Hernandez_Andrea_Espino_Heriberto.tex
\usepackage[discipline=computer-science]{template/liverpoolthesis}
\usepackage[utf8]{inputenc}
\usepackage[T1]{fontenc}
\usepackage[spanish,es-nodecimaldot]{babel}

% Las referencias internas y las citas usan el mismo azul que los títulos.
\hypersetup{
  linkcolor = LiverpoolBlue,
  citecolor = LiverpoolBlue,
  urlcolor = LiverpoolBlue
}

% La marca de cada capítulo aparece en el encabezado, arriba de la regla.
\addto\captionsspanish{\renewcommand{\chaptername}{Ejercicio}}
\providecommand{\figureautorefname}{Figura}
\addto\captionsspanish{\renewcommand{\figureautorefname}{Figura}}

\graphicspath{{figuras/}}
\setlength{\headheight}{28pt}

% Inserta una captura de código y/o de ejecución.
% El tercer argumento es la etiqueta; el argumento opcional controla la altura.
\newcommand{\incluirevidencia}[4][0.70\textheight]{%
  \begin{figure}[H]
    \centering
    \includegraphics[
      width=0.95\textwidth,
      height=#1,
      keepaspectratio
    ]{#2}
    \caption{#3}
    \label{#4}
  \end{figure}%
}

% ---------- Datos de la portada ----------
\title{Práctica de laboratorio 1: Creación de procesos e hilos en C y Java}
\author{%
  Heriberto Espino Montelongo\\
  ID: 175199\\[1em]
  Andrea Hernandez Huerta\\
  ID: 175523\\[1em]
  \text{Profesor:} Dr. Victor Giovanni Morales Murillo}
\faculty{Universidad de las Américas Puebla (UDLAP)}
\school{Licenciatura en Ciencia de Datos}
\degree{Curso: Programación Concurrente}
\date{Fecha de entrega: \today}
\subject{Creación y ejecución de procesos e hilos en C y Java}

\graphicspath{{figuras/}{../latex/figuras/}}

\title{Actividad 9: Semáforo binario y contador}
\author{%
  Andrea Hernandez Huerta\\
  ID: 175523\\[1em]
  Heriberto Espino Montelongo\\
  ID: 175199\\[1em]
  \text{Profesor:} Dr. Victor Giovanni Morales Murillo}
\faculty{Universidad de las Américas Puebla (UDLAP)}
\school{Licenciatura en Ciencia de Datos}
\degree{Curso: Programación Concurrente}
\date{Fecha de entrega: \today}
\subject{Semáforo binario y semáforo contador}

\addto\captionsspanish{\renewcommand{\chaptername}{Parte}}
\providecommand{\figureautorefname}{Figura}
\addto\captionsspanish{\renewcommand{\figureautorefname}{Figura}}

\begin{document}

\frontmatter
\maketitle
\tableofcontents
\listoffigures
\listoftables

\mainmatter

\chapter{Semáforo binario: cuenta bancaria}

\textbf{Instrucción.} Crear una cuenta bancaria compartida con saldo inicial de
\$1000 y cinco hilos que realicen depósitos y retiros. Utilizar un semáforo
binario inicializado en 1. Cada hilo debe solicitar el permiso, consultar y
modificar el saldo y liberar el permiso. Se debe comprobar que solamente un
hilo modifique el saldo a la vez y realizar una prueba de escritorio con dos
hilos.

\section{Entradas}

La cuenta bancaria inicia con un saldo de \$1000. Los cinco hilos comparten la
misma instancia de \texttt{CuentaBancaria} y la misma instancia de
\texttt{BinarySemaphore}. El semáforo se crea con \texttt{true}, equivalente
a un único permiso disponible.

Las operaciones utilizadas en la ejecución son tres depósitos de \$200,
\$300 y \$50, y dos retiros de \$150 y \$100. Por ello, el saldo final esperado
es
\[
1000+200-150+300-100+50=1300.
\]

\section{Procesamiento}

La clase \texttt{BinarySemaphore} mantiene una variable booleana
\texttt{value}. Cuando su valor es \texttt{true}, el permiso está disponible;
cuando es \texttt{false}, el permiso está ocupado.

La operación \texttt{P()} es un método \texttt{synchronized}. Mientras el
permiso esté ocupado, el hilo ejecuta \texttt{Util.mywait(this)}. La espera
libera temporalmente el monitor para permitir que otro hilo ejecute
\texttt{V()}. Al obtener el permiso, \texttt{P()} cambia el valor a
\texttt{false}.

La operación \texttt{V()} también es \texttt{synchronized}: cambia el valor a
\texttt{true} y ejecuta \texttt{notify()} para despertar a uno de los hilos
bloqueados.

La clase \texttt{CuentaBancaria} no utiliza \texttt{synchronized}; el acceso
al saldo queda protegido externamente por el semáforo. Cada
\texttt{OperacionBancaria} ejecuta \texttt{P()} antes de entrar a la sección
crítica y \texttt{V()} en la sección de salida.

\section{Estructura de clases}

\begin{table}[H]
\centering
\caption{Responsabilidad de las clases de la parte 1.}
\label{tab:clases-binario}
\begin{tabularx}{\textwidth}{P{3.7cm} L P{2.6cm}}
\toprule
\textbf{Clase} & \textbf{Responsabilidad} & \textbf{Uso de synchronized}\\
\midrule
Principal & Crea la cuenta, el semáforo y los cinco hilos; utiliza
\texttt{start()} y \texttt{join()}. & No lo requiere.\\
\texttt{BinarySemaphore} & Controla el único permiso mediante
\texttt{P()} y \texttt{V()}. & \texttt{P()} y \texttt{V()} son
\texttt{synchronized}.\\
\texttt{CuentaBancaria} & Almacena el saldo y realiza depósitos y retiros. &
No es necesario si el acceso está protegido por el semáforo.\\
\texttt{OperacionBancaria} & Define la sección crítica ejecutada por cada
hilo y utiliza \texttt{P()} antes y \texttt{V()} después de modificar la
cuenta. & No lo requiere.\\
\bottomrule
\end{tabularx}
\end{table}

\section{Utilidad compartida}

La clase \texttt{Util} concentra las esperas bloqueadas y las pausas usadas
por ambos ejercicios. Su método \texttt{mywait()} ejecuta \texttt{wait()} sobre
el monitor recibido y conserva la interrupción del hilo; \texttt{mysleep()}
permite simular el tiempo de una operación sin convertir la espera en un ciclo
activo.

\incluirevidencia{Util.png}
  {Clase Java de utilidades compartidas por ambos ejercicios.}
  {fig:util}

\section{Conceptos utilizados}

El semáforo binario realiza exclusión mutua porque existe un solo permiso. Si
un hilo ya lo obtuvo, los demás se bloquean en \texttt{P()} hasta que el hilo
que se encuentra en la sección crítica ejecute \texttt{V()}. A diferencia de
la espera activa estudiada con Test-and-Set, la operación de espera utiliza
\texttt{wait()}, por lo que el hilo bloqueado libera el monitor y no permanece
consultando continuamente el valor del semáforo.

En las diapositivas, \texttt{P()} corresponde a solicitar un permiso, también
llamada \texttt{acquire()}, \texttt{wait()} o \texttt{down()}.
\texttt{V()} corresponde a liberar un permiso, también llamada
\texttt{release()}, \texttt{signal()} o \texttt{up()}.

\section{Prueba de escritorio}

Para la prueba se consideran dos hilos. El hilo A deposita \$200 y el hilo B
retira \$150. El semáforo inicia disponible y el saldo inicia en \$1000.

\begin{table}[H]
\centering
\caption{Prueba de escritorio del semáforo binario con dos hilos.}
\label{tab:prueba-binario}
\begin{tabularx}{\textwidth}{c L c c L}
\toprule
\textbf{Paso} & \textbf{Acción} & \textbf{Semáforo} & \textbf{Saldo} &
\textbf{Resultado}\\
\midrule
0 & Estado inicial & 1 & 1000 & Permiso disponible.\\
1 & Hilo A ejecuta \texttt{P()} & 0 & 1000 & A entra a la sección crítica.\\
2 & Hilo B ejecuta \texttt{P()} & 0 & 1000 & B se bloquea porque el permiso está ocupado.\\
3 & A deposita \$200 & 0 & 1200 & Solo A modifica la cuenta.\\
4 & A ejecuta \texttt{V()} & 1 & 1200 & Se libera el permiso y B es despertado.\\
5 & B completa \texttt{P()} & 0 & 1200 & B obtiene el permiso.\\
6 & B retira \$150 & 0 & 1050 & Solo B modifica la cuenta.\\
7 & B ejecuta \texttt{V()} & 1 & 1050 & El permiso queda disponible.\\
\bottomrule
\end{tabularx}
\end{table}

\section{Salida esperada y comprobación}

El programa muestra cuándo cada hilo entra a la sección crítica, el saldo antes
y después de la operación y el saldo final. Las cinco operaciones producen un
saldo esperado de \$1300. La comprobación adicional debe reportar un máximo de
un hilo modificando el saldo al mismo tiempo.

\incluirevidencia{resultados/1.png}
  {Ejecución del semáforo binario: saldo final de \$1300 y máximo de un hilo
   modificando el saldo simultáneamente.}
  {fig:resultado-binario}

\section{Evidencia de la implementación}

\incluirevidencia{1/binario_codigo_01.png}
  {Implementación de \texttt{BinarySemaphore}, \texttt{P()} y \texttt{V()}.}
  {fig:binario-codigo-01}

\incluirevidencia[0.55\textwidth]{1/binario_codigo_02.png}
  {Implementación de \texttt{OperacionBancaria} y sus datos compartidos.}
  {fig:binario-codigo-02}
\vspace{-.5cm}
\incluirevidencia[0.95\textwidth]{1/binario_codigo_03.png}
  {Sección crítica de la operación bancaria y actualización del saldo.}
  {fig:binario-codigo-03}

\incluirevidencia{1/binario_codigo_04.png}
  {Creación de los cinco hilos y comprobación del semáforo binario.}
  {fig:binario-codigo-04}

\chapter{Semáforo contador: impresoras}

\textbf{Instrucción.} Simular un laboratorio con tres impresoras y ocho hilos
que desean imprimir. Utilizar un semáforo contador inicializado en 3. Cada hilo
debe solicitar una impresora, indicar que comienza a imprimir, esperar un
tiempo para simular la impresión, indicar que termina y liberar la impresora.
Se debe verificar que nunca existan más de tres hilos imprimiendo
simultáneamente y realizar una prueba de escritorio con dos impresoras y cuatro
hilos.

\section{Entradas}

El recurso está formado por tres impresoras equivalentes. El programa crea un
semáforo contador con valor inicial 3 y ocho hilos que comparten el mismo
semáforo. Cada tarea mantiene ocupada una impresora durante un intervalo breve
para hacer visible la concurrencia en la salida.

\section{Procesamiento}

La clase \texttt{CountingSemaphore} mantiene un entero \texttt{value}, como en
el pseudocódigo mostrado en clase. En \texttt{P()} se decrementa el valor. Si
el resultado es negativo, el hilo ejecuta una espera bloqueada. Un valor
negativo representa que existen hilos esperando un recurso.

En \texttt{V()} se incrementa el valor. Si después del incremento el valor es
menor o igual que cero, existe al menos un hilo esperando y se ejecuta
\texttt{notify()} para despertar a uno. Tanto \texttt{P()} como \texttt{V()}
son métodos \texttt{synchronized}.

Cada hilo solicita una impresora con \texttt{P()}, realiza la simulación de
impresión y finalmente devuelve el permiso con \texttt{V()}. La pausa usada
para simular la impresión no constituye el mecanismo de sincronización;
únicamente representa el tiempo durante el cual el recurso permanece ocupado.

\section{Estructura de clases}

\begin{table}[H]
\centering
\caption{Responsabilidad de las clases de la parte 2.}
\label{tab:clases-contador}
\begin{tabularx}{\textwidth}{P{3.8cm} L P{2.7cm}}
\toprule
\textbf{Clase} & \textbf{Responsabilidad} & \textbf{Uso de synchronized}\\
\midrule
Principal & Crea el semáforo con tres permisos y los ocho hilos; utiliza
\texttt{start()} y \texttt{join()}. & No lo requiere.\\
\texttt{CountingSemaphore} & Mantiene el número de permisos y controla la
espera mediante \texttt{P()} y \texttt{V()}. & \texttt{P()} y \texttt{V()}
son \texttt{synchronized}.\\
\texttt{TareaImpresion} & Solicita una impresora, simula su utilización y la
libera al terminar. & No lo requiere.\\
\bottomrule
\end{tabularx}
\end{table}

\section{Prueba de escritorio}

La prueba utiliza dos impresoras y cuatro hilos A, B, C y D. Los dos primeros
hilos obtienen un recurso. Los siguientes disminuyen el contador por debajo de
cero y quedan bloqueados. Cuando A y B liberan sus impresoras, C y D son
despertados de uno en uno.

\begin{table}[H]
\centering
\caption{Prueba de escritorio del semáforo contador con dos impresoras y cuatro hilos.}
\label{tab:prueba-contador}
\begin{tabularx}{\textwidth}{c L c P{2.0cm} L}
\toprule
\textbf{Paso} & \textbf{Acción} & \textbf{Semáforo} & \textbf{Cola} &
\textbf{Resultado}\\
\midrule
0 & Estado inicial & 2 & Vacía & Dos impresoras disponibles.\\
1 & Hilo A ejecuta \texttt{P()} & 1 & Vacía & A utiliza una impresora.\\
2 & Hilo B ejecuta \texttt{P()} & 0 & Vacía & B utiliza la otra impresora.\\
3 & Hilo C ejecuta \texttt{P()} & -1 & C & C se bloquea.\\
4 & Hilo D ejecuta \texttt{P()} & -2 & C, D & D se bloquea.\\
5 & Hilo A ejecuta \texttt{V()} & -1 & D & C despierta y utiliza la impresora liberada.\\
6 & Hilo B ejecuta \texttt{V()} & 0 & Vacía & D despierta y utiliza la impresora liberada.\\
7 & Hilo C ejecuta \texttt{V()} & 1 & Vacía & Queda una impresora disponible.\\
8 & Hilo D ejecuta \texttt{V()} & 2 & Vacía & Quedan dos impresoras disponibles.\\
\bottomrule
\end{tabularx}
\end{table}

\section{Salida esperada y comprobación}

La salida identifica cuándo inicia y termina cada hilo y muestra cuántos hilos
se encuentran imprimiendo. El orden exacto puede cambiar entre ejecuciones
porque depende de la planificación de los hilos. La propiedad que sí debe
mantenerse es que el máximo observado nunca sea mayor que 3.

\incluirevidencia{resultados/2.png}
  {Ejecución del semáforo contador: máximo de tres hilos imprimiendo
   simultáneamente, igual al límite permitido.}
  {fig:resultado-contador}

\section{Evidencia de la implementación}

\incluirevidencia{2/contador_codigo_01.png}
  {Implementación de \texttt{CountingSemaphore}, \texttt{P()} y \texttt{V()}.}
  {fig:contador-codigo-01}

\incluirevidencia{2/contador_codigo_02.png}
  {Implementación de \texttt{TareaImpresion} y control del máximo simultáneo.}
  {fig:contador-codigo-02}

\incluirevidencia{2/contador_codigo_03.png}
  {Creación de las ocho tareas y ejecución del semáforo contador.}
  {fig:contador-codigo-03}

\chapter{Comparación y conclusión}

\section{Comparación}

\begin{table}[H]
\centering
\caption{Comparación entre semáforo binario y semáforo contador.}
\label{tab:comparacion}
\begin{tabularx}{\textwidth}{P{4.0cm} C C}
\toprule
\textbf{Aspecto} & \textbf{Semáforo binario} & \textbf{Semáforo contador}\\
\midrule
Valor inicial en la actividad & 1 (\texttt{true}) & 3\\
Permisos disponibles inicialmente & 1 & 3\\
Recurso controlado & Cuenta bancaria & Tres impresoras\\
Número máximo de hilos simultáneos & 1 & 3\\
Cuándo se bloquea un hilo & Cuando el único permiso está ocupado &
Cuando no queda un permiso disponible\\
\bottomrule
\end{tabularx}
\end{table}

\section{Conclusión}

Los dos semáforos coordinan hilos mediante permisos, pero representan problemas
diferentes. El semáforo binario mantiene un solo permiso y por ello se utiliza
para exclusión mutua sobre un recurso compartido, como el saldo de la cuenta.
El semáforo contador mantiene varios permisos y permite que una cantidad
limitada de hilos utilice simultáneamente recursos equivalentes, como las tres
impresoras.

En ambos casos, \texttt{P()} solicita un permiso y puede bloquear al hilo,
mientras que \texttt{V()} devuelve un permiso y puede despertar a un hilo que
espera. La diferencia principal está en la cantidad de permisos que representa
el semáforo. Las pruebas de escritorio muestran esta diferencia: el semáforo
binario obliga a que los accesos a la cuenta ocurran uno por uno, mientras que
el semáforo contador permite hasta tres impresiones concurrentes sin superar el
número de recursos disponibles.

\cleardoublepage
\phantomsection
\addcontentsline{toc}{chapter}{Referencias}
\nocite{*}
\bibliography{template/references}

\end{document}
